% =====================================================================
% Chapter 7 — Route A Results
% =====================================================================
\chapter{Route~A Results}
\label{ch:resultsA}

This chapter reports the experimental results of \routeA{}. It presents the
registration outcomes for all five \suop{} categories
(\secref{sec:resultsA_registration}), the mesh statistics
(\secref{sec:resultsA_meshing}), the variant results (view-aware and ray TSDF,
\secref{sec:resultsA_variants}) and a discussion of what the numbers do and do
not prove (\secref{sec:resultsA_interpretation}). The full pipeline ran on all
five object categories with the 24-view curation described in
\secref{ch:dataset}; every run produced fused clouds and both meshes, i.e.\
5/5 categories completed the full geometric pipeline.

\section{Registration results}
\label{sec:resultsA_registration}

\tabref{tab:registration_results} summarises the pose-graph stage for the five
objects. All $24 \times 5$ views were connected: for every category the
maximum-fitness spanning tree reached 24/24 views and no segment was left
outside the graph. The number of reliable pairwise edges (fitness $\geq 0.25$)
varies between 131 (drum) and 197 (dummy) of the 276 pairs tested, and the
number of loop closures --- edges beyond the 23 MST edges --- ranges from 108
to 174. The high loop-closure counts are important: they mean the pose graph
was strongly over-constrained, so drift along the tree was corrected by many
independent constraints rather than trusted as a chain.

\begin{table}[H]
  \centering
  \small
  \begin{adjustbox}{max width=\textwidth}
  \begin{tabular}{@{}lcccccc@{}}
    \toprule
    Object & Views & Reliable edges & MST edges & Loop closures & Fused points (raw) & Fused points (7.5~mm) \\
    \midrule
    Tyre  & 24 & 172 & 23 & 149 & 37{,}031 & 33{,}970 \\
    Dummy & 24 & 197 & 23 & 174 & 22{,}692 & 21{,}421 \\
    Drum  & 24 & 131 & 23 & 108 & 106{,}436 & 99{,}544 \\
    Chair & 24 & 189 & 23 & 166 & 30{,}519 & 26{,}027 \\
    Net   & 24 & 162 & 23 & 139 & 31{,}288 & 28{,}251 \\
    \bottomrule
  \end{tabular}
  \end{adjustbox}
  \caption{Route~A pose-graph results. 276 pairwise registrations were tested
  per object ($\binom{24}{2}$); edges with fitness $\geq 0.25$ entered the
  graph; the maximum-fitness spanning tree used 23 edges and every remaining
  reliable edge became a loop closure. Fused counts are the pose-corrected
  points before and after the 7.5~mm fusion voxel. All 24 views were reachable
  for every object.}
  \label{tab:registration_results}
\end{table}

The pairwise computation cost was substantial but feasible: between 912~s
(dummy) and 1098~s (chair) per object on the CPU for the 276 pairwise
FPFH/RANSAC/ICP evaluations, dominated by the RANSAC global alignment. This is
an offline cost and was accepted; the per-view ICP refinement itself is
milliseconds once initialised.

\begin{figure}[H]
  \centering
  \includegraphics[width=0.92\textwidth]{fig_chair_segments.png}
  \caption{Segmented chair views used as \routeA{} input (one panel per
  accepted case). Each panel shows the isolated chair cluster after seabed-plane
  removal and scale-consistent clustering at the published $0.40$~m diagonal
  scale. The views are partial and noisy --- note the variable leg coverage ---
  yet registration must still align them into one frame.}
  \label{fig:chair_segments}
\end{figure}

\section{Meshing results}
\label{sec:resultsA_meshing}

Both Screened Poisson (depth 8) and BPA (four radii) produced clean meshes for
all five objects. \tabref{tab:meshing_results} reports vertex and triangle
counts after trimming and cleaning (largest connected component only).

\begin{table}[H]
  \centering
  \small
  \begin{adjustbox}{max width=\textwidth}
  \begin{tabular}{@{}lcccccc@{}}
    \toprule
    Object & Fused points & Poisson verts & Poisson tris & BPA verts & BPA tris \\
    \midrule
    Tyre  & 33{,}841 & 69{,}819 & 139{,}684 & 14{,}143 & 24{,}009 \\
    Dummy & 21{,}306 & 29{,}787 & 59{,}419 & 11{,}194 & 19{,}331 \\
    Drum  & 99{,}678 & 125{,}598 & 251{,}677 & 33{,}429 & 53{,}649 \\
    Chair & 25{,}872 & 85{,}891 & 172{,}484 & 9{,}534 & 15{,}317 \\
    Net   & 28{,}240 & 76{,}319 & 152{,}721 & 5{,}220 & 8{,}313 \\
    \bottomrule
  \end{tabular}
  \end{adjustbox}
  \caption{Route~A mesh statistics. Poisson (depth 8, trimmed at the 2\%
  lowest-density vertices) produces smooth closed surfaces with many more
  vertices; BPA produces sparser surfaces that preserve genuine gaps. The
  divergence between the two column pairs is itself the evidence of
  measurement uncertainty (see \secref{sec:resultsA_interpretation}).}
  \label{tab:meshing_results}
\end{table}

\begin{figure}[H]
  \centering
  \includegraphics[width=0.94\textwidth]{fig_chair_reconstruction.png}
  \caption{Route~A chair: fused measured cloud and the two meshes. The chair's
  open frame is a stress test: leg-like structures are registered across many
  views, and the difference between the closed Poisson surface and the open BPA
  surface is clearly visible in the seat underside and between legs.}
  \label{fig:chair_reconstruction}
\end{figure}

\begin{figure}[H]
  \centering
  \includegraphics[width=0.94\textwidth]{fig_tyre_reconstruction.png}
  \caption{Route~A tyre: fused cloud, Poisson and BPA surfaces. The toroidal
  rim is recovered; the inner wall and far side, which no 3~m view observed
  fully, show the Poisson-vs-BPA disagreement most clearly.}
  \label{fig:tyre_reconstruction}
\end{figure}

\begin{figure}[H]
  \centering
  \includegraphics[width=0.94\textwidth]{fig_dummy_reconstruction.png}
  \caption{Route~A dummy: fused cloud, Poisson and BPA surfaces. The vertical
  mannequin shape and the shoulder/arm junctions are visible; the head and
  torso are the best-supported regions.}
  \label{fig:dummy_reconstruction}
\end{figure}

\subsection{Turntable views}
\figref{fig:chair_turntable} shows the chair reconstruction rotated around the
vertical axis. The turntable renders make the 3D consistency of the fused
model auditable from all sides rather than from a single favourable angle.

\begin{figure}[H]
  \centering
  \includegraphics[width=0.98\textwidth]{fig_chair_turntable.png}
  \caption{Chair reconstruction turntable: the same fused model viewed from
  eight azimuths. Consistency across views is the visual equivalent of the
  graph's loop-closure redundancy.}
  \label{fig:chair_turntable}
\end{figure}

\begin{figure}[H]
  \centering
  \includegraphics[width=0.98\textwidth]{fig_dummy_turntable.png}
  \caption{Dummy reconstruction turntable. The rear and side views show where
  the measured coverage is thinnest and where Poisson and BPA therefore
  diverge most.}
  \label{fig:dummy_turntable}
\end{figure}

\section{Variant results}
\label{sec:resultsA_variants}

The two variants of \secref{sec:routea_variants} were run on the chair and
tyre, the two categories with the strongest segmentation. Their metrics are
reported in \tabref{tab:routeA_metrics}. The metric
\emph{observed-to-surface} measures, for every accepted registered return, its
distance to the reconstructed surface; a small median means the surface is
faithful to the measurements. \emph{Surface supported} is the fraction of
surface points with a measured return within a support radius (40~mm for the
view-aware variant, the support limit for the TSDF).

\begin{table}[H]
  \centering
  \small
  \begin{adjustbox}{max width=\textwidth}
  \begin{tabular}{@{}lcccccc@{}}
    \toprule
    Variant & Object & Accepted cases & Fused returns & Median (mm) & P95 (mm) & Supported \\
    \midrule
    View-aware SE(2) & Chair & 7 & 972 & 4.14 & 19.31 & 0.63 ($\le 20$~mm) \\
    View-aware SE(2) & Tyre  & 3 & 5{,}000 & 2.91 & 15.82 & 1.00 ($\le 40$~mm) \\
    Ray TSDF & Chair & 2 & 2{,}205 & 2.92 & 9.89 & 1.00 \\
    Ray TSDF & Tyre  & 5 & 7{,}982 & 4.99 & 20.05 & 1.00 \\
    \bottomrule
  \end{tabular}
  \end{adjustbox}
  \caption{Surface-fit metrics for the two \routeA{} variants. Median
  observed-to-surface residuals of $2.9$--$5.0$~mm confirm that the fused
  surfaces sit tightly on the registered returns; the P95 values show that a
  small tail of returns (often multipath or mis-segmented points) lies a few
  centimetres off. ``Supported'' is the fraction of surface within the stated
  radius of a measured return.}
  \label{tab:routeA_metrics}
\end{table}

\begin{itemize}
  \item \textbf{View-aware SE(2).} The chair accepted 7 of its candidate
  cases (972 fused returns) with a median residual of 4.14~mm; the tyre
  accepted 3 (5{,}000 returns) with median 2.91~mm and essentially total
  surface support within 40~mm. The small accepted-case counts are the honest
  cost of the explicit acceptance gate: cases whose planar fit is weak are
  rejected rather than silently fused.
  \item \textbf{Ray TSDF.} The sonar-native variant fused 2,205 chair returns
  from 2 views and 7,982 tyre returns from 5 views into 4--5~mm voxel grids.
  Its surfaces are the most conservative: unsupported vertices are removed, so
  the fraction of multiview-supported surface is only 0.176 (chair) and 0.499
  (tyre) --- the TSDF reports \emph{per-vertex} how many views support it,
  which is exactly the coverage information a sonar reconstruction should
  provide.
\end{itemize}

\subsection{Model-assisted hybrid outputs}
For completeness, the project also produced ``recognizable'' hybrid outputs in
which a dimensionally constrained template (folding chair, parametric tyre)
completes the registered evidence. These are \emph{model-assisted}, not
measurement-derived: the chair completion places inferred frame members
between measured joints, and the tyre completion uses the measured width with a
parametric profile. The tyre hybrid's larger residuals (median 18.5~mm,
P95 65.8~mm) quantify the cost of forcing a template onto sparse evidence, and
make clear why the purely measured variants are preferred for claims of
geometric accuracy.

\section{Interpretation: what these numbers do and do not prove}
\label{sec:resultsA_interpretation}

The numbers above establish four things rigorously:

\begin{enumerate}
  \item \textbf{Pipeline completion.} All five categories completed the full
  registration$\to$fusion$\to$meshing chain; the graphs are fully connected
  (24/24 views), so no object was abandoned.
  \item \textbf{Measurement consistency.} Median observed-to-surface residuals
  of a few millimetres mean that, wherever the sonar did measure, the fused
  surface agrees with the measurements at a scale comparable to the
  segmentation and registration tolerances.
  \item \textbf{Multi-view benefit.} The large loop-closure counts and the
  visual agreement across turntable views show that the fused model is
  self-consistent, which a single scan cannot be.
  \item \textbf{Uncertainty localisation.} The systematic Poisson-vs-BPA
  divergence in the least-supported regions (chair underside, tyre inner wall,
  dummy rear) is a map of where the reconstruction is extrapolated rather than
  measured.
\end{enumerate}

What the numbers do \emph{not} prove is absolute geometric accuracy: there is
no external CAD or scanned reference in \suop{}, so ``4~mm residual'' is a
self-consistency figure, not an error against truth. Likewise, mesh vertex
counts are indicators of pipeline behaviour, not quality. These caveats are
central to the discussion in \secref{ch:discussion} and motivate the
calibrated-acquisition future work of \secref{ch:conclusions}.
